\section{Magisterial, Canonical, and Liturgical Corpus}
\label{app:freemasonry-corpus}

\subsection{How to read the inventory}

The first table is complete under the inclusion rule stated at the beginning of this article: it contains the located universal papal or Roman act whose principal object is Freemasonry, whose named adjacent society is expressly absorbed into the developing Masonic/secret-society discipline, whose canon directly governs Masonic enrollment, or whose Roman interpretation controls the present question. An act can be historically superseded as penal law while remaining essential evidence of development. “Present role” therefore never means that an abrogated censure is still applied.

The later tables contain verified adjacent sources that materially frame the controversy but are broader in subject, and the official liturgical or disciplinary sources needed to classify exorcism prayers. They are intentionally not called a complete list of every papal allusion to secret societies or every private deliverance formula. Full texts or official archival witnesses are linked in the references. The summaries identify operative content; they are not substitute promulgations.

\subsection{Controlling and materially direct corpus}

\begin{fourstable}{Date and authority}{Act and genre}{Material contribution}{Present role}
28 Apr. 1738, Clement XII & \emph{In eminenti apostolatus specula}, apostolic constitution & First universal papal prohibition: interreligious lodge association, oath-bound secrecy, danger to Church and public order; reserved excommunication. & Historical origin of continuous prohibition; text officially reproduced by \emph{Quo graviora}.\\
18 May 1751, Benedict XIV & \emph{Providas Romanorum Pontificum}, constitution & Confirms that Clement's prohibition remains, restates secrecy, oath, religious mixture, and public danger, and rejects attempted evasion. & Historical confirmation; officially reproduced by \emph{Quo graviora}.\\
13 Sept. 1821, Pius VII & \emph{Ecclesiam a Jesu Christo}, bull/constitution on the Carbonari & Condemns the named Carbonari society and situates new clandestine bodies within the danger addressed by predecessors. & Direct for Carbonari and material to the papal genealogy; not proof that Carbonari and every lodge are identical.\\
13 Mar. 1825, Leo XII & \emph{Quo graviora}, bull & Reproduces the three earlier acts, confirms them, and addresses old and newly named clandestine associations. & Official early corpus witness and renewed discipline.\\
25 Sept. 1865, Pius IX & \emph{Multiplices inter}, consistorial allocution & Direct exposition against the Masonic sect; recounts predecessors, rejects the claim of harmlessness, renews condemnation and censures. & Major pre-Leonine synthesis; empirical claims remain historically situated.\\
12 Oct. 1869, Pius IX & \emph{Apostolicae Sedis moderationi}, apostolic constitution, reserved-censure list no.~IV & Reorganizes \emph{latae sententiae} censures and names Masonic, Carbonari, and like sects plotting openly or secretly against Church or legitimate powers, including material support and non-denunciation of hidden leaders. & Superseded penal law; documentary bridge to 1917 canon 2335.\\
20 Apr. 1884, Leo XIII & \emph{Humanum genus}, encyclical on Freemasonry & Systematic diagnosis of naturalism, indifferentism, secrecy, social program, and remedy; confirms predecessors; explicitly distinguishes generic judgment from individual members. & Principal magisterial exposition; read with later teaching and paragraph 11's qualification.\\
15 Oct. 1890, Leo XIII & \emph{Dall'alto dell'Apostolico Seggio}, encyclical/letter on Freemasonry in Italy & Applies the analysis to Italian anti-clerical laws, schools, property, public institutions, and the papacy; cites public Masonic programs. & Direct historical application, not a universal description of every national body.\\
8 Dec. 1892, Leo XIII & \emph{Inimica vis}, encyclical to the Italian bishops & Reiterates the Italian conflict, naturalist program, pastoral vigilance, doctrine, Catholic association, and prayer. & Direct episcopal instruction in its historical setting.\\
8 Dec. 1892, Leo XIII & \emph{Custodi di quella fede}, encyclical/letter to the Italian people & Popular companion to \emph{Inimica vis}; calls Catholics to guard faith, formation, press, association, charity, and prayer against Masonic influence. & Direct lay-pastoral application; source of the incompatibility formula quoted by the CDF in 1985.\\
27 May 1917 / 19 May 1918, Benedict XV & Promulgation and effective date of the 1917 CIC; canon 2335 & Names the Masonic sect and like associations plotting against Church or legitimate powers; imposes \emph{ipso facto} excommunication reserved to the Apostolic See. & Abrogated by the 1983 Code; must not be applied as current law.\\
19 July 1974 / 17 Feb. 1981, CDF & Restricted interpretive letter; public declaration in AAS 73 (1981) 240--241 & The 1974 letter allowed strict penal interpretation for bodies actually plotting; widespread overreading prompted the 1981 declaration that discipline remained and conferences lacked power for a general derogation. & Explains the transition and prevents use of the 1974 letter as a general permission.\\
25 Jan. / 27 Nov. 1983, John Paul II & Promulgation and effective date of the 1983 CIC; canon 1374 & Uses the broad category “association which plots against the Church”; provides a just penalty for joining and interdict for promotion or office. & Current Latin penal category, revised in wording within Book VI in 2021; no named automatic Masonic excommunication.\\
26 Nov. 1983, John Paul II / CDF & \emph{Declaration on Masonic Associations} & Negative judgment unchanged; principles irreconcilable; membership forbidden; enrolled faithful in grave sin and not to receive Communion; no local derogation. & Controlling universal doctrinal and sacramental judgment.\\
23 Feb. 1985, CDF & \emph{Reflections on the Declaration}, doctrinal explanatory article & Locates incompatibility below practical hostility: progressive symbol-system, relativizing ritual frame, two relations to God, faith reduced to one institutional expression; explains objective grave sin. & Authoritative dicasterial explanation of rationale and scope, distinct in genre from the 1983 declaration.\\
13 Nov. 2023, Francis / DDF & Response to Bishop Julito Cortes, note approved \emph{ex Audientia} & Reaffirms forbidden active membership and 1983 declaration; addresses formal, knowing enrollment and embrace of principles; applies to clerics; urges coordinated catechesis. & Latest located universal Roman application through the as-of date; Philippine occasion, universal doctrinal norm.\\
\end{fourstable}

\subsection{Verified adjacent papal sources}

\begin{threestable}{Source}{Direct subject}{Why it belongs in the interpretive frame}
Pius VIII, \emph{Traditi humilitati}, 24 May 1829 & Broad programmatic encyclical on contemporary religious and social errors. & Named by \emph{Humanum genus} among predecessor interventions; illuminates indifferentism and clandestine opposition without becoming a Masonry-only act.\\
Gregory XVI, \emph{Mirari vos}, 15 Aug. 1832 & Indifferentism, conscience and press claims, Church-state order, and associations in the post-revolutionary crisis. & Also named in \emph{Humanum genus}; supplies the nineteenth-century doctrinal-political context.\\
Pius IX, \emph{Qui pluribus}, 9 Nov. 1846 & Rationalism, indifferentism, Revelation, and enemies of Church and society. & Cited by \emph{Humanum genus} among Pius IX's relevant acts; not represented here as solely about Masonry.\\
Pius IX, \emph{Quanta cura}, 8 Dec. 1864 & Errors concerning religion, Church, conscience, and civil order. & Frames the doctrinal and political world immediately preceding \emph{Multiplices inter}; exact claims must be read in genre and context.\\
Leo XIII, \emph{Officio sanctissimo}, 22 Dec. 1887, especially no.~12 & The Church in Bavaria; includes a direct warning against Masonic influence, especially upon youth. & A materially direct paragraph within a broader encyclical, following \emph{Humanum genus}.\\
Pius XI, \emph{Caritate Christi compulsi}, 3 May 1932 & Prayer, economic crisis, atheistic movements, and organized enemies of religion. & Shows the wider secret-society and anti-religious frame; does not replace the specific Masonic corpus.\\
Second Vatican Council, \emph{Dignitatis humanae}, 7 Dec. 1965 & Civil religious freedom grounded in human dignity and bounded by public order. & Necessary later doctrinal development: civil immunity from coercion is not religious indifferentism or ritual equivalence.\\
Francis, \emph{Pascite gregem Dei}, 23 May 2021 & Promulgation of revised Book VI of the 1983 Code, effective 8 Dec. 2021. & Fixes the current edition and effective date of canon 1374 used in this article.\\
\end{threestable}

\subsection{Verified adjacent exorcism and prayer-discipline sources}

\begin{fourstable}{Date and authority}{Text}{Material rule or finding}{Relation to Masonry}
18 May 1890, Leo XIII & \emph{Exorcismus in Satanam et angelos apostaticos}, ASS 23 & Official formula against Satan and fallen angels; rescript concerns bishops and priests legitimately authorized by ordinaries. & Contains \latin{omnis congregatio et secta diabolica} but never names Freemasonry; adjacent to, not part of, the direct Masonic corpus.\\
29 Sept. 1985, CDF & Letter to ordinaries on the norms governing exorcism & Reserves exorcism according to CIC can.~1172; bars lay use of the Leonine formula and unauthorized interrogation of demons. & Governs attempts to use the Leonine text or private formulas as self-administered anti-Masonic rites.\\
26 Jan. 1999, Holy See presentation & Revised \emph{De Exorcismis et Supplicationibus Quibusdam} & Major rite, prayers concerning places or things and attacks on the Church, and private supplications are kept distinct. & No Masonry-specific formula was located in the official public inventory consulted.\\
14 Sept. 2000, CDF & Instruction on prayers for healing & Requires dependence upon the diocesan bishop; keeps exorcism distinct from healing services and out of Mass, sacraments, and the Hours. & Prevents a Masonry-renunciation service from being presented as liturgy or inserted into sacramental celebration.\\
1 Nov. 2024, Spanish Episcopal Conference & Doctrinal note on intergenerational healing & Rejects inherited personal guilt and a special-ritual theory of diseases caused by ancestors' unforgiven sins; affirms real natural and social consequences. & Particular episcopal application that directly tests automatic ancestral-curse claims in modern Masonry prayers.\\
\end{fourstable}

\subsection{Continuity and change in one view}

\begin{center}
\begin{tikzpicture}[node distance=5mm]
\node[guide node, text width=.19\linewidth] (early) {1738--1869\\secrecy, oaths, religious mixture, hostile societies; named censures};
\node[guide node, right=of early, text width=.19\linewidth] (leo) {1884--1892\\naturalism and indifferentism systematized; positive Catholic response};
\node[guide node, right=of leo, text width=.19\linewidth] (codes) {1917--1983\\named automatic censure gives way to broader penal category};
\node[guide emphasis, right=of codes, text width=.19\linewidth] (current) {1983--2026\\principial incompatibility, grave sin, Communion discipline, no local derogation};
\draw[guide arrow] (early) -- (leo);
\draw[guide arrow] (leo) -- (codes);
\draw[guide arrow] (codes) -- (current);
\end{tikzpicture}
\end{center}

The continuity is the prohibition and its doctrinal core. The change is the penal technique, the specification of the deeper rationale, the Church's developed account of civil religious freedom, and the pastoral setting. One cannot turn legal development into reversal; one cannot turn doctrinal continuity into the claim that every old censure or political formulation remains current law unchanged.
